% GENERATED FILE. Edit canonical lessons, then run npm run generate:books.
% canonical-source-hash: fnv1a64:b33999043a1b5bcc
% canonical-lessons: SA-C376-yogyah, SA-C376-samarthah, SA-C376-niyacchati, SA-C376-niyantranam, SA-C376-asrayate

\chapter{Who Can Do It}
\label{ch:sa-who-can-do-it}
\hlchaptermodality{376}

\begin{chapteropening}
\textbf{By the end of this chapter:} \emph{I can say who is fit for a task and who keeps it under control.}

Complete the last lesson of chapter 376: I can say who is fit for a task and who keeps it under control.
\end{chapteropening}

\section[\texorpdfstring{\sk{योग्यः} (yogyaḥ)}{योग्यः (yogyaḥ)}]{\sk{योग्यः} (yogyaḥ) — qualified; fit}

\label{lesson:SA-C376-yogyah}



\begin{quote}
Before the new one: say the Sanskrit for a drought, then the Sanskrit for remedies.
\end{quote}

\subsection*{You'll want to know: \sk{योग्यः}}
\textbf{\sk{योग्यः}} — \emph{yogyaḥ} — \textquotedblleft{}qualified; fit\textquotedblright{}.

\begin{quote}
\sk{यदि} \sk{अरुणः} \sk{योग्यः} \sk{अस्ति}, \sk{तर्हि} \sk{सः} \sk{नायकः} \sk{भवेत्।}
\end{quote}

\begin{quote}
\emph{yadi aruṇaḥ yogyaḥ asti, tarhi saḥ nāyakaḥ bhavet.}
\end{quote}

\begin{quote}
If Arun is qualified, he should be the leader.
\end{quote}

\begin{cousinweb}
\textbf{\sk{योग्यः}} (\emph{yogyaḥ}) is built on \emph{yoga}, \textquotedblleft{}a yoking, union\textquotedblright{}, from the root √yuj, \textquotedblleft{}to join\textquotedblright{}; MW gives \emph{yogya} \textquotedblleft{}fit for the yoke\textquotedblright{} and \textquotedblleft{}useful, proper, qualified, able\textquotedblright{} (MW).
\end{cousinweb}

\begin{grammarlens}[title={What you've built}]
One more word for saying what happens if.
\end{grammarlens}

\subsection*{Your turn}
\begin{itemize}
\raggedright
  \item \emph{Say it:} \emph{yogyaḥ}
  \item \emph{Say it:} \emph{yogyaḥ}, once more
  \item \emph{Say it:} \emph{yadi aruṇaḥ yogyaḥ asti, tarhi saḥ nāyakaḥ bhavet.} — If Arun is qualified, he should be the leader.
  \item \emph{Recall:} say the Sanskrit for a drought, then the Sanskrit for remedies, then say \emph{yogyaḥ} again
\end{itemize}

\begin{tcolorbox}[breakable,colback=teal!4,colframe=teal!35!black,title={Before you move on}]
What does \sk{योग्यः} mean? (\textquotedblleft{}Qualified; fit\textquotedblright{}.) Say it once more.
\end{tcolorbox}
\section[\texorpdfstring{\sk{समर्थः} (samarthaḥ)}{समर्थः (samarthaḥ)}]{\sk{समर्थः} (samarthaḥ) — capable; able}

\label{lesson:SA-C376-samarthah}



\begin{quote}
Before the new one: say the Sanskrit for remedies, then the Sanskrit for qualified.
\end{quote}

\subsection*{You'll want to know: \sk{समर्थः}}
\textbf{\sk{समर्थः}} — \emph{samarthaḥ} — \textquotedblleft{}capable; able\textquotedblright{}.

\begin{quote}
\sk{यदि} \sk{सः} \sk{समर्थः}, \sk{तर्हि} \sk{कार्यं} \sk{करिष्यति।}
\end{quote}

\begin{quote}
\emph{yadi saḥ samarthaḥ, tarhi kāryaṃ kariṣyati.}
\end{quote}

\begin{quote}
If he is capable, he will do the work.
\end{quote}

\begin{cousinweb}
\textbf{\sk{समर्थः}} (\emph{samarthaḥ}) is built from \emph{sam}, \textquotedblleft{}together\textquotedblright{}, and \emph{artha}, \textquotedblleft{}aim\textquotedblright{}; MW gives \emph{sam-artha} \textquotedblleft{}having a suitable aim\textquotedblright{} and \textquotedblleft{}competent, capable, able\textquotedblright{} (MW).
\end{cousinweb}

\begin{grammarlens}[title={What you've built}]
One more word for saying what happens if.
\end{grammarlens}

\subsection*{Your turn}
\begin{itemize}
\raggedright
  \item \emph{Say it:} \emph{samarthaḥ}
  \item \emph{Say it:} \emph{samarthaḥ}, once more
  \item \emph{Say it:} \emph{yadi saḥ samarthaḥ, tarhi kāryaṃ kariṣyati.} — If he is capable, he will do the work.
  \item \emph{Recall:} say the Sanskrit for remedies, then the Sanskrit for qualified, then say \emph{samarthaḥ} again
\end{itemize}

\begin{tcolorbox}[breakable,colback=teal!4,colframe=teal!35!black,title={Before you move on}]
What does \sk{समर्थः} mean? (\textquotedblleft{}Capable; able\textquotedblright{}.) Say it once more.
\end{tcolorbox}
\section[\texorpdfstring{\sk{नियच्छति} (niyacchati)}{नियच्छति (niyacchati)}]{\sk{नियच्छति} (niyacchati) — controls; restrains}

\label{lesson:SA-C376-niyacchati}



\begin{quote}
Before the new one: say the Sanskrit for qualified, then the Sanskrit for capable.
\end{quote}

\subsection*{You'll want to know: \sk{नियच्छति}}
\textbf{\sk{नियच्छति}} — \emph{niyacchati} — \textquotedblleft{}controls; restrains\textquotedblright{}.

\begin{quote}
\sk{यदि} \sk{भवान्} \sk{मनः} \sk{नियच्छति}, \sk{तर्हि} \sk{शान्तिः} \sk{भविष्यति।}
\end{quote}

\begin{quote}
\emph{yadi bhavān manaḥ niyacchati, tarhi śāntiḥ bhaviṣyati.}
\end{quote}

\begin{quote}
If you control your mind, there will be peace.
\end{quote}

Past (he …): \sk{सः} \sk{क्रोधं} \sk{नियतवान्}, \emph{saḥ krodhaṃ niyatavān}.

\begin{cousinweb}
\textbf{\sk{नियच्छति}} (\emph{niyacchati}) is built on \emph{ni}, \textquotedblleft{}down\textquotedblright{}, and the root √yam, \textquotedblleft{}to hold, to restrain\textquotedblright{}; MW gives \emph{ni-√yam} \textquotedblleft{}to restrain, to curb, to control, to govern\textquotedblright{} (MW).
\end{cousinweb}

\begin{grammarlens}[title={What you've built}]
One more word for saying what happens if.
\end{grammarlens}

\subsection*{Your turn}
\begin{itemize}
\raggedright
  \item \emph{Say it:} \emph{niyacchati}
  \item \emph{Say it:} \emph{niyacchati}, once more
  \item \emph{Say it:} \emph{yadi bhavān manaḥ niyacchati, tarhi śāntiḥ bhaviṣyati.} — If you control your mind, there will be peace.
  \item \emph{Recall:} say the Sanskrit for qualified, then the Sanskrit for capable, then say \emph{niyacchati} again
\end{itemize}

\begin{tcolorbox}[breakable,colback=teal!4,colframe=teal!35!black,title={Before you move on}]
What does \sk{नियच्छति} mean? (\textquotedblleft{}Controls; restrains\textquotedblright{}.) Say it once more.
\end{tcolorbox}
\section[\texorpdfstring{\sk{नियन्त्रणम्} (niyantraṇam)}{नियन्त्रणम् (niyantraṇam)}]{\sk{नियन्त्रणम्} (niyantraṇam) — control; restraint}

\label{lesson:SA-C376-niyantranam}



\begin{quote}
Before the new one: say the Sanskrit for capable, then the Sanskrit for controls.
\end{quote}

\subsection*{You'll want to know: \sk{नियन्त्रणम्}}
\textbf{\sk{नियन्त्रणम्}} — \emph{niyantraṇam} — \textquotedblleft{}control; restraint\textquotedblright{}.

\begin{quote}
\sk{यदि} \sk{चालकस्य} \sk{नियन्त्रणं} \sk{न} \sk{अस्ति}, \sk{तर्हि} \sk{भयम्} \sk{अस्ति।}
\end{quote}

\begin{quote}
\emph{yadi cālakasya niyantraṇaṃ na asti, tarhi bhayam asti.}
\end{quote}

\begin{quote}
If the driver has no control, there is danger.
\end{quote}

A neuter noun.

\begin{cousinweb}
\textbf{\sk{नियन्त्रणम्}} (\emph{niyantraṇam}) is built on \emph{ni} and the root √yantr, \textquotedblleft{}to restrain, to bind\textquotedblright{}; MW gives \emph{ni-yantraṇa} \textquotedblleft{}restraining, checking, governing\textquotedblright{} (MW).
\end{cousinweb}

\begin{grammarlens}[title={What you've built}]
One more word for saying what happens if.
\end{grammarlens}

\subsection*{Your turn}
\begin{itemize}
\raggedright
  \item \emph{Say it:} \emph{niyantraṇam}
  \item \emph{Say it:} \emph{niyantraṇam}, once more
  \item \emph{Say it:} \emph{yadi cālakasya niyantraṇaṃ na asti, tarhi bhayam asti.} — If the driver has no control, there is danger.
  \item \emph{Recall:} say the Sanskrit for capable, then the Sanskrit for controls, then say \emph{niyantraṇam} again
\end{itemize}

\begin{tcolorbox}[breakable,colback=teal!4,colframe=teal!35!black,title={Before you move on}]
What does \sk{नियन्त्रणम्} mean? (\textquotedblleft{}Control; restraint\textquotedblright{}.) Say it once more.
\end{tcolorbox}
\section[\texorpdfstring{\sk{आश्रयते} (āśrayate)}{आश्रयते (āśrayate)}]{\sk{आश्रयते} (āśrayate) — takes refuge in; leans on}

\label{lesson:SA-C376-asrayate}



\begin{quote}
Before the new one: say the Sanskrit for controls, then the Sanskrit for control.
\end{quote}

\subsection*{You'll want to know: \sk{आश्रयते}}
\textbf{\sk{आश्रयते}} — \emph{āśrayate} — \textquotedblleft{}takes refuge in; leans on\textquotedblright{}.

\begin{quote}
\sk{यदि} \sk{संकटम्} \sk{अस्ति}, \sk{तर्हि} \sk{बालः} \sk{मातरम्} \sk{आश्रयते।}
\end{quote}

\begin{quote}
\emph{yadi saṅkaṭam asti, tarhi bālaḥ mātaram āśrayate.}
\end{quote}

\begin{quote}
If there is danger, the child takes refuge in its mother.
\end{quote}

Past (he …): \sk{सः} \sk{मित्रम्} \sk{आश्रितवान्}, \emph{saḥ mitram āśritavān}.

\begin{cousinweb}
\textbf{\sk{आश्रयते}} (\emph{āśrayate}) is built on \emph{ā} and the root √śri, \textquotedblleft{}to lean on\textquotedblright{}; MW gives \emph{ā-√śri} \textquotedblleft{}to seek refuge in\textquotedblright{} and \textquotedblleft{}to depend on\textquotedblright{} (MW).
\end{cousinweb}

\begin{grammarlens}[title={What you've built}]
One more word for saying what happens if.
\end{grammarlens}

\subsection*{Your turn}
\begin{itemize}
\raggedright
  \item \emph{Say it:} \emph{āśrayate}
  \item \emph{Say it:} \emph{āśrayate}, once more
  \item \emph{Say it:} \emph{yadi saṅkaṭam asti, tarhi bālaḥ mātaram āśrayate.} — If there is danger, the child takes refuge in its mother.
  \item \emph{Recall:} say the Sanskrit for controls, then the Sanskrit for control, then say \emph{āśrayate} again
\end{itemize}

\begin{tcolorbox}[breakable,colback=teal!4,colframe=teal!35!black,title={Before you move on}]
What does \sk{आश्रयते} mean? (\textquotedblleft{}Takes refuge in; leans on\textquotedblright{}.) Say it once more.
\end{tcolorbox}
